\documentclass{beamer}

\usepackage[T2A]{fontenc}
\usepackage[utf8]{inputenc}
\usepackage[english,russian]{babel}

% \usepackage{mathpazo}      % шрифт Palatino + математика
\usefonttheme{serif}       % beamer использует шрифт с засечками

\usetheme{CambridgeUS}
\usecolortheme{wolverine}

\setbeamertemplate{caption}[numbered]

% \usepackage{pgfpages}
% \setbeameroption{show notes on second screen}

\title[Математическое моделирование]{Основные принципы построения математических моделей}
\author[Гркикян М.Э.]{Гркикян Мисак Эдикович}
\date{09.10.2026}

% \setbeamertemplate{title page}{
% 	\centering
% 	\vskip1em
% 	{\usebeamerfont{title}\usebeamercolor[fg]{title}\inserttitle\par}
% 	\vskip1em
% 	{\usebeamerfont{author}\insertauthor\par}
% }

\begin{document}

	\begin{frame}
		\titlepage
	\end{frame}

    \begin{frame}{1. Применение фундаментальных законов природы}
        \textbf{Постановка задачи.} Оценить заряд батареи ноутбука в конце дня.

        \medskip
        \textbf{Обозначения.} Пусть $E(t)$ - заряд батареи (Вт$\cdot$ч) в момент времени $t$;
        $\Delta E_{\text{зар}}$ - энергия, полученная при зарядке;
        $\Delta E_{\text{расх}}$ - энергия, потраченная на работу.

        \medskip
        \textbf{Модель} (закон сохранения энергии):
        \[
            E(0) + \Delta E_{\text{зар}} = E(t) + \Delta E_{\text{расх}}.
        \]

        \medskip
        \textbf{Пример.} $E(0) = 50$ Вт$\cdot$ч,
        $\Delta E_{\text{зар}} = 30$ Вт$\cdot$ч,
        $\Delta E_{\text{расх}} = 60$ Вт$\cdot$ч. Тогда
        \[
            E(t) = 50 + 30 - 60 = 20\ \text{Вт}\cdot\text{ч}.
        \]

        \textbf{Ограничения модели.} Не учитывается изнашивание аккумулятора.
        Потеря энергии, когда нагревается. Когда устройство садится быстрее на морозе, хотя энергии столько же.
    \end{frame}

    \begin{frame}{2. Применение вариационных принципов}
        \textbf{Постановка задачи.} Есть два экзамена. Нужно так распределить время
        на подготовку, чтобы потратить его как можно меньше, но набрать
        нужный балл по каждому экзамену.

        \medskip
        \textbf{Обозначения.} Пусть $x$ и $y$ - часы потраченные на подготовку на первый и второй экзамен; 
        $a_1$ и $a_2$ - сколько баллов ты «зарабатываешь» за один час подготовки к первому и второму экзаменам;
        $B_1$ и $B_2$ - минимальный балл, который тебе нужен по первому и второму экзаменам.

    \end{frame}

    \begin{frame}{2. Применение вариационных принципов}
        \textbf{Вариационный принцип:}
        \[
            L = x + y \to \min,
        \]
        \[
            \begin{cases}
                x \ge 0,\ y \ge 0, \\
                a_1 x \ge B_1, \\
                a_2 y \ge B_2.
            \end{cases}
        \]
        
        \textbf{Ограничения модели.}
        \begin{itemize}
            \item Усталость — после многих часов эффективность падает.
            \item Разная сложность билетов.
        \end{itemize}
    \end{frame}

    \begin{frame}{3. Применение аналогий}
        \textbf{Ключевая мысль.} Две разные задачи могут описываться одной и той же
        математической моделью - тогда решение одной автоматически даёт
        решение другой.

        \medskip
        \textbf{Постановка задачи.} За день нужно успеть и кардио, и силовую
        тренировку. Хочется потратить как можно меньше суммарного времени,
        но при этом набрать минимальную норму по каждому виду нагрузки.

        \medskip
        \textbf{Обозначения.} Пусть $x$ и $y$ - минуты, потраченные на кардио
        и силовую; $a_1$ и $a_2$ - «эффект» одной минуты кардио
        и силовой (условные единицы нагрузки в минуту); $B_1$ и $B_2$ -
        минимальная норма нагрузки по кардио и силовой (эффект).
    \end{frame}

    \begin{frame}{3. Применение аналогий}
        \textbf{Вариационный принцип:}
        \[
            L = x + y \to \min,
        \]
        \[
            \begin{cases}
                x \ge 0,\ y \ge 0, \\
                a_1 x \ge B_1, \\
                a_2 y \ge B_2.
            \end{cases}
        \]

        \medskip
        \textbf{Ограничения модели.} Не учитывается усталость (эффективность
        падает со временем) и взаимное влияние тренировок друг на друга.
    \end{frame}

    \begin{frame}{4. Иерархический подход}
        \textbf{Постановка задачи.} Дан отсортированный массив из $n$ чисел.
        Нужно найти в нём заданное число $x$. Алгоритм - бинарный поиск.

        \medskip
        \textbf{Идея.} На каждом шаге массив делится пополам, и поиск
        продолжается только в той половине, где может быть $x$.
        Каждый уровень работает с вдвое меньшим массивом - это и есть иерархия.

        \medskip
        \textbf{Уровни:}
        \begin{itemize}
            \item уровень 1 - весь массив длины $n$;
            \item уровень 2 - половина длины $n/2$;
            \item уровень 3 - четверть длины $n/4$;
            \item \dots
            \item уровень $k$ - отрезок длины $n/2^k$.
        \end{itemize}
    \end{frame}

    \begin{frame}{4. Иерархический подход}
        \textbf{Число действий.} Пусть $T(n)$ - сколько шагов нужно,
        чтобы найти элемент в массиве длины $n$. Тогда
        \[
            T(n) = T\!\left(\frac{n}{2}\right) + 1,
        \]
        где $T(n/2)$ - работа на следующем уровне (массив вдвое короче),
        а $+1$ - одно сравнение с серединой на текущем уровне.

        \medskip
        Разворачивая рекурсию:
        \[
            T(n) = T\!\left(\frac{n}{2^k}\right) + k.
        \]
        Рекурсия останавливается при $n/2^k = 1$, откуда $k = \log_2 n$. 
    \end{frame}

    \begin{frame}{4. Иерархический подход}
        \textbf{Числовой пример.} Массив из 8 чисел:
        \[
            [\,1,\ 3,\ 5,\ 7,\ 9,\ 11,\ 13,\ 15\,], \quad x = 13.
        \]
        \begin{itemize}
            \item уровень 1: середина 7 $\Rightarrow$ $13 > 7$, идём вправо;
            \item уровень 2: $[9, 11, 13, 15]$, середина 11 $\Rightarrow$ вправо;
            \item уровень 3: $[13, 15]$, середина 13 $\Rightarrow$ нашли.
        \end{itemize}
        Всего 3 шага вместо 8 при переборе: $\log_2 8 = 3$.

        \medskip
        \textbf{Ограничения модели.}
        \begin{itemize}
            \item Массив должен быть отсортирован.
        \end{itemize}
    \end{frame}
\end{document}